\documentclass[11pt,a4paper]{article}
\usepackage[margin=2.3cm,headheight=15pt]{geometry}
\usepackage[T1]{fontenc}
\usepackage[utf8]{inputenc}
\usepackage[spanish,es-nodecimaldot]{babel}
\usepackage{lmodern,microtype}
\usepackage{graphicx,float,booktabs,tabularx,array,xcolor}
\usepackage{amsmath,amssymb}
\usepackage{caption,subcaption,fancyhdr,titlesec}
\usepackage{hyperref}
\usepackage{minted}
\usepackage{csquotes}
\usepackage[backend=biber,style=ieee]{biblatex}
\addbibresource{references.bib}

\definecolor{navy}{RGB}{0,32,91}
\definecolor{pale}{RGB}{241,244,248}
\hypersetup{colorlinks=true,linkcolor=navy,citecolor=navy,urlcolor=navy,
  pdftitle={Cifrados clásicos: Alberti y rejilla giratoria}}
\pagestyle{fancy}\fancyhf{}
\fancyhead[L]{\small Cifrados clásicos}
\fancyhead[R]{\small Seguridad en computación}
\fancyfoot[C]{\thepage}
\setlength{\parskip}{0.45em}
\setlength{\parindent}{0pt}
\titleformat{\section}{\Large\bfseries\color{navy}}{\thesection}{0.7em}{}
\titleformat{\subsection}{\large\bfseries\color{navy}}{\thesubsection}{0.7em}{}
\captionsetup{font=small,labelfont=bf}

\newcommand{\capturaarchivo}[2]{%
  \begin{minipage}[t]{0.48\textwidth}
    \centering\includegraphics[width=\linewidth,height=4.7cm,keepaspectratio]{#1}\\[0.3em]
    {\small #2}
  \end{minipage}}

\begin{document}
\begin{titlepage}
  \centering
  \vspace*{0.5cm}
  \includegraphics[width=0.22\textwidth]{csunsa.pdf}\par\vspace{1.1cm}
  {\Large UNIVERSIDAD NACIONAL DE SAN AGUSTÍN\par}
  \vspace{0.25cm}
  {\large Facultad de Ingeniería de Producción y Servicios\\
  Escuela Profesional de Ciencia de la Computación\par}
  \vspace{1.4cm}
  {\Large\bfseries SEGURIDAD EN COMPUTACIÓN\par}
  \vspace{0.8cm}
  {\huge\bfseries Cifrados clásicos:\\[0.2em] disco de Alberti y rejilla giratoria\par}
  \vspace{1.2cm}
  {\large Práctica: implementación de cifrado y descifrado\par}
  \vspace{0.55cm}
  {\small Aplicación: \href{https://elquesohelado.github.io/Alberti_Fleissner/}{elquesohelado.github.io/Alberti\_Fleissner}\\
  Repositorio: \href{https://github.com/ElQuesoHelado/Alberti_Fleissner}{github.com/ElQuesoHelado/Alberti\_Fleissner}\par}
  \vfill
  \begin{tabular}{rl}
    Integrantes: & Alva Cornejo, Jose Javier\\
                 & Carpio Paiva, Cesar Gonzalo\\
                 & Palomino Rivadeneyra, Misael Jesús\\
                 & Zapana Pariapaza, Gonzalo Rodrigo
  \end{tabular}
  \vfill
  {Arequipa -- Perú\\\today\par}
\end{titlepage}

\tableofcontents
\newpage

\section{Objetivo y alcance}
Este trabajo describe y presenta una implementación web educativa de dos cifrados clásicos: la sustitución polialfabética inspirada en el disco de Alberti y la transposición mediante rejilla giratoria de Fleissner. Para cada método se explica el principio, el alfabeto o clave, el proceso inverso de descifrado, las limitaciones de la implementación y espacios reservados para evidencias de ejecución.

La aplicación se ejecuta en el navegador. Está implementada en TypeScript con React y Vite; las funciones de cifrado y descifrado se encuentran en \texttt{src/ciphers.ts}. No requiere servidor para procesar los mensajes. Estos algoritmos se presentan con fines académicos y no deben usarse para proteger información sensible.

\section{Preprocesado y convenciones compartidas}
Antes de cifrar, el perfil usado en este trabajo convierte el texto a mayúsculas, elimina diacríticos (\texttt{á} pasa a \texttt{A}; \texttt{ñ} pasa a \texttt{N}), elimina espacios y saltos de línea, y filtra la entrada para conservar únicamente las letras latinas A--Z y los dígitos 0--9. Así se eliminan signos y puntuación. La aplicación permite cambiar estas opciones, pero para reproducir los resultados se debe usar el perfil descrito y compartir la configuración JSON. La vista previa permite comprobar el texto resultante.

Por ejemplo, \texttt{Ataque al amanecer!} se normaliza como \texttt{ATAQUEALAMANECER}; el texto descifrado aparece unido y en mayúsculas.

El objeto JSON que acompaña al criptograma guarda las opciones de preprocesado. Quien descifra debe recibir tanto el criptograma como esa configuración y seleccionar el algoritmo correspondiente. El texto original no se reconstruye exactamente si el preprocesado eliminó diferencias de caja, tildes o espacios; el resultado es el texto normalizado. En la rejilla, además, se guarda la longitud para retirar el relleno. En Alberti, los caracteres ajenos a los discos se conservan o eliminan según la opción seleccionada.

\section{Cifrado y descifrado con el disco de Alberti}
\subsection{Principio del algoritmo}
El método sustituye cada símbolo claro por un símbolo de un alfabeto cifrante. La implementación usa un anillo exterior como alfabeto de entrada y un alfabeto interior como alfabeto de salida. El giro inicial y un incremento periódico determinan qué posición del alfabeto interior se usa. El principio del disco móvil fue descrito por Leon Battista Alberti en \textit{De cifris} (siglo XV) \cite{alberti}.

Sean $P_i$ la posición del símbolo claro en el anillo exterior, $m$ la longitud de los alfabetos, $o$ el giro inicial, $a$ el incremento y $q$ el período. Para el símbolo cifrable número $i$, la posición cifrante es
\[
  C_i = (P_i + o + \lfloor i/q\rfloor a)\bmod m.
\]
La numeración de $i$ solo avanza al procesar un símbolo presente en el disco. El descifrado busca el símbolo en el alfabeto interior y aplica la operación inversa:
\[
  P_i = (C_i - o - \lfloor i/q\rfloor a)\bmod m.
\]

\subsection{Alfabeto, clave y limitaciones}
Los alfabetos predeterminados tienen 24 símbolos y deben contener la misma cantidad de símbolos únicos:
\begin{itemize}
  \item Exterior (claro): \texttt{ABCDEFGILMNOPQRSTVXZ1234}.
  \item Interior (cifrado): \texttt{usqomkhfdbacegilnprtxz\&y}.
\end{itemize}
La clave de descifrado contiene ambos alfabetos, el giro inicial $o$, el incremento $a$ y el período $q$. En la interfaz, el período debe ser al menos 1. Los alfabetos son configurables y el validador exige igual longitud y ausencia de duplicados.

La variante implementada no es una réplica completa de todas las convenciones históricas del disco. No incluye indicadores externos para anunciar cambios arbitrarios de giro. Las letras que no aparecen en el alfabeto exterior (por ejemplo, J, H, K, U, W y Ñ antes de normalizar) se consideran no soportadas: se conservan sin cifrar o se eliminan según la opción elegida. Si se conservan símbolos ajenos, esos símbolos no aportan al contador periódico. La sustitución clásica no proporciona seguridad criptográfica moderna.

\subsection{Procedimiento de cifrado}
\begin{enumerate}
  \item Aplicar las opciones de preprocesado al texto claro.
  \item Validar que ambos alfabetos tengan longitud igual y símbolos únicos.
  \item Para cada carácter reconocido, hallar su posición exterior, calcular el giro según $o$, $a$ y $q$, y emitir el símbolo interior correspondiente.
  \item Aplicar la política configurada a los caracteres que no estén en el alfabeto.
  \item Entregar el criptograma y el JSON de parámetros.
\end{enumerate}

\subsection{Procedimiento de descifrado}
\begin{enumerate}
  \item Recibir el criptograma y el JSON producido durante el cifrado.
  \item Validar los alfabetos y recuperar giro, incremento, período y tratamiento de caracteres.
  \item Contar los símbolos reconocidos en el mismo orden; buscar cada uno en el alfabeto interior y restar el giro que le corresponde módulo $m$.
  \item Mostrar el texto claro normalizado. No se restauran tildes, espacios ni caja que hayan sido descartados antes del cifrado.
\end{enumerate}

\subsection{Evidencias de ejecución}
La primera captura muestra el cifrado y su descifrado; la segunda muestra la configuración del disco.
\begin{figure}[H]\centering
  \capturaarchivo{alberti_1.png}{Cifrado y descifrado}\hfill\capturaarchivo{alberti_2.png}{Configuración}
  \caption{Resultados y parámetros del disco de Alberti.}\label{fig:alberti-evidencias}
\end{figure}

\section{Cifrado y descifrado con la rejilla giratoria}
\subsection{Principio del algoritmo}
La rejilla giratoria es un cifrado de transposición: cambia el orden de los caracteres, pero no sustituye sus valores. La configuración usada en este trabajo es una rejilla de 4 × 4, girada en sentido horario en cuatro posiciones. Las perforaciones se leen por filas en cada giro. Una máscara válida visita cada celda del tablero una sola vez \cite{fleissner1881,histocrypt2018}.

La máscara se representa con \texttt{X} para una celda perforada y \texttt{.} para una celda cerrada:
\begin{center}\texttt{X... / ...X / ..X. / .X..}\end{center}
Durante el cifrado se escribe el texto a través de las perforaciones en cada orientación y se lee el tablero completo por filas. En el descifrado se coloca el criptograma en el mismo orden por filas y se recuperan los caracteres recorriendo las perforaciones en el mismo orden de giros. La configuración compartida contiene estas posiciones, el sentido y el número de giros, además de la longitud original del texto.

\subsection{Alfabeto, clave y limitaciones}
El alfabeto de entrada es \texttt{ABCDEFGHIJKLMNOPQRSTUVWXYZ0123456789}: letras latinas mayúsculas sin Ñ y dígitos. El preprocesado convierte a mayúsculas, quita tildes, convierte Ñ en N y elimina espacios, signos y cualquier otro símbolo fuera del alfabeto. El equipo de descifrado recibe la longitud original junto con el criptograma; la configuración JSON de la aplicación también registra la longitud y los parámetros de la rejilla.

La rejilla indicada tiene una perforación por cada órbita y completa 16 celdas tras cuatro giros. Si el texto ocupa menos celdas, las posiciones restantes se rellenan con \texttt{X} hasta completar el bloque. La longitud original se comunica junto con el cifrado y permite retirar el relleno al descifrar. La pérdida de caracteres o de esta longitud puede impedir la recuperación. La transposición no oculta las frecuencias de los caracteres y no ofrece seguridad moderna.

\subsection{Procedimiento de cifrado}
\begin{enumerate}
  \item Aplicar el preprocesado y el filtro de caracteres seleccionados.
  \item Validar la rejilla: cada órbita de rotación debe tener exactamente una perforación.
  \item Recorrer por filas los huecos de cada una de las cuatro posiciones horarias y llenar las celdas.
  \item Completar las celdas vacías del bloque con \texttt{X}; leer el tablero por filas y compartir la longitud original y los parámetros.
\end{enumerate}

\subsection{Procedimiento de descifrado}
\begin{enumerate}
  \item Cargar el criptograma junto con el JSON de configuración.
  \item Validar tamaño, perforaciones y que la longitud del criptograma coincida con los bloques esperados.
  \item Recorrer las mismas perforaciones y giros, extraer hasta la longitud registrada y descartar el relleno.
  \item Presentar el texto filtrado y normalizado. Los espacios y símbolos que se descartaron no pueden reconstruirse.
\end{enumerate}

\subsection{Evidencias de ejecución}
La primera captura muestra el cifrado y su descifrado; la segunda muestra el preprocesado y la configuración de la rejilla, incluida su longitud original.
\begin{figure}[H]\centering
  \capturaarchivo{rejilla_1.png}{Cifrado y descifrado}\hfill\capturaarchivo{rejilla_2.png}{Preprocesado y configuración}
  \caption{Resultados y parámetros de la rejilla giratoria.}\label{fig:grille-evidencias}
\end{figure}

\section{Implementación y código}
La aplicación está escrita en TypeScript y usa React para la interfaz y Vite para el desarrollo y empaquetado web. El siguiente listado incluye las funciones implementadas de preprocesado, cifrado y descifrado de ambos métodos, tomadas directamente de \texttt{src/ciphers.ts} del proyecto \cite{proyecto}.

\subsection{Código fuente de los algoritmos}
\inputminted[fontsize=\scriptsize,breaklines,linenos,frame=single,framesep=2mm]{typescript}{../src/ciphers.ts}

\section{Conclusiones}
Se implementaron dos principios criptográficos distintos: sustitución con alfabetos móviles para Alberti y transposición con una máscara rotatoria para Fleissner. El preprocesado y la configuración compartida permiten fijar convenciones compatibles entre quien cifra y quien descifra. La recuperación depende de usar los mismos alfabetos o la misma rejilla, parámetros y longitud, y el texto se recupera en forma normalizada cuando el preprocesado elimina información. Ambos métodos son apropiados para estudiar algoritmos clásicos, pero no para proteger datos reales.

\printbibliography
\end{document}
